\begin{abstract}
	电机作为工业生产中最重要的动力设备之一，其安全可靠运行对保障生产效率和设备安全具有重要意义。然而，电机在长期运行过程中不可避免地会出现各种故障，若不能及时发现和处理，将导致停机事故和经济损失。传统的电机故障诊断方法主要依赖人工经验判断和信号处理技术，存在诊断准确率低、泛化能力差等问题。

	本文针对电机故障诊断中特征提取困难、诊断精度不足的问题，提出了一种基于深度学习的电机故障诊断方法。主要研究内容包括：（1）分析了电机常见故障类型及其振动信号特征，建立了电机故障振动信号数据集；（2）提出了一种基于卷积神经网络的端到端故障诊断方法，直接从原始振动信号中自动学习故障特征，避免了传统方法中复杂的人工特征提取过程；（3）引入迁移学习策略，解决了工业场景中小样本条件下模型训练困难的问题；（4）设计了一种多尺度特征融合网络结构，有效提高了模型对不同类型故障的识别能力。

	实验结果表明，本文提出的方法在电机故障诊断任务中取得了95.6\%的平均准确率，优于传统的支持向量机和随机森林等方法，同时在小样本场景下仍能保持较高的诊断精度，验证了该方法的有效性和实用性。

	\keywords{电机故障诊断，深度学习，卷积神经网络，迁移学习，特征提取}
\end{abstract}

\begin{enabstract}
	As one of the most important power equipment in industrial production, the safe and reliable operation of motors is of great significance to ensuring production efficiency and equipment safety. However, various faults inevitably occur during long-term operation of motors. If not detected and handled in time, they will lead to shutdown accidents and economic losses. Traditional motor fault diagnosis methods mainly rely on manual experience judgment and signal processing technology, which have problems such as low diagnostic accuracy and poor generalization ability.

	This thesis proposes a motor fault diagnosis method based on deep learning to address the problems of difficult feature extraction and insufficient diagnostic accuracy in motor fault diagnosis. The main research contents include: (1) analyzing common fault types of motors and their vibration signal characteristics, and establishing a motor fault vibration signal dataset; (2) proposing an end-to-end fault diagnosis method based on convolutional neural network, which automatically learns fault features directly from raw vibration signals, avoiding the complex manual feature extraction process in traditional methods; (3) introducing transfer learning strategy to solve the problem of difficult model training under small sample conditions in industrial scenarios; (4) designing a multi-scale feature fusion network structure to effectively improve the model's ability to recognize different types of faults.

	Experimental results show that the method proposed in this thesis achieves an average accuracy of 95.6\% in motor fault diagnosis tasks, outperforming traditional methods such as support vector machine and random forest. Meanwhile, it can still maintain high diagnostic accuracy under small sample scenarios, verifying the effectiveness and practicality of the proposed method.

	\enkeywords{Motor Fault Diagnosis, Deep Learning, Convolutional Neural Network, Transfer Learning, Feature Extraction}
\end{enabstract}